\documentclass[12pt,a4paper]{article}
\usepackage[polish]{babel}
\usepackage[T1]{fontenc}
\usepackage[utf8x]{inputenc}
\usepackage{hyperref}
\usepackage{url}
\usepackage[]{algorithm2e}
\usepackage{listings}
\usepackage{amssymb}
\usepackage{amsmath}
\usepackage{float}

\usepackage{color}
\usepackage{listings}

\lstloadlanguages{% Check Dokumentation for further languages ...
	C,
	C++,
	csh,
	Java
}

\definecolor{red}{rgb}{0.6,0,0} % for strings
\definecolor{blue}{rgb}{0,0,0.6}
\definecolor{green}{rgb}{0,0.8,0}
\definecolor{cyan}{rgb}{0.0,0.6,0.6}

\lstset{
	language=csh,
	basicstyle=\footnotesize\ttfamily,
	numbers=left,
	numberstyle=\tiny,
	numbersep=5pt,
	tabsize=2,
	extendedchars=true,
	breaklines=true,
	frame=b,
	stringstyle=\color{blue}\ttfamily,
	showspaces=false,
	showtabs=false,
	xleftmargin=17pt,
	framexleftmargin=17pt,
	framexrightmargin=5pt,
	framexbottommargin=4pt,
	commentstyle=\color{green},
	morecomment=[l]{//}, %use comment-line-style!
	morecomment=[s]{/*}{*/}, %for multiline comments
	showstringspaces=false,
	morekeywords={ abstract, event, new, struct,
		as, explicit, null, switch,
		base, extern, object, this,
		bool, false, operator, throw,
		break, finally, out, true,
		byte, fixed, override, try,
		case, float, params, typeof,
		catch, for, private, uint,
		char, foreach, protected, ulong,
		checked, goto, public, unchecked,
		class, if, readonly, unsafe,
		const, implicit, ref, ushort,
		continue, in, return, using,
		decimal, int, sbyte, virtual,
		default, interface, sealed, volatile,
		delegate, internal, short, void,
		do, is, sizeof, while,
		double, lock, stackalloc,
		else, long, static,
		enum, namespace, string},
	keywordstyle=\color{cyan},
	identifierstyle=\color{red},
}
\usepackage{caption}
\DeclareCaptionFont{white}{\color{white}}
\DeclareCaptionFormat{listing}{\colorbox{blue}{\parbox{\textwidth}{\hspace{15pt}#1#2#3}}}
\captionsetup[lstlisting]{format=listing,labelfont=white,textfont=white, singlelinecheck=false, margin=0pt, font={bf,footnotesize}}


\addtolength{\hoffset}{-1.5cm}
\addtolength{\marginparwidth}{-1.5cm}
\addtolength{\textwidth}{3cm}
\addtolength{\voffset}{-1cm}
\addtolength{\textheight}{2.5cm}
\setlength{\topmargin}{0cm}
\setlength{\headheight}{0cm}

\begin{document}

\title{Symulacja wybuchu epidemii zombie\\\small{Dokumentacja projektu systemu wieloagentowego}}
\author{Rafal}
\date{\today}

\maketitle
\tableofcontents
\newpage

% ========== CZĘŚĆ I: OPIS SYSTEMU ==========
\section{Model Systemu Informatycznego}

\subsection{Definicja Modelu}
System ``Zombie Outbreak Simulation'' jest symulacją wieloagentową opracowywaną
w środowisku gry 2D. Model systemu stanowi szachownica o wymiarach 128x128
komórek, na której poruszają się pięć typów autonomicznych agentów
posiadających różne cele i zachowania. Każdy agent podejmuje decyzje w oparciu
o obserwacje otoczenia i stanu wewnętrznego, stanowiąc część zamodelowanego
ekosystemu zmagan ludzi zombie.

\subsection{Uzasadnienie w Kontekście Budowanego Systemu}
Model wieloagentowy (MAS - Multi-Agent System) pozwala na symulowanie złożonych
interakcji między jednostkami o różnych celach strategicznych. W kontekście
symulacji epidemiologicznej, taki system umożliwia:
\begin{itemize}
	\item Badanie dynamiki rozprzestrzeniania się infekcji
	\item Analiza efektywności strategii defensywnych (medycy, żołnierze)
	\item Obserwacja emergentnych zachowań wynikających z prostych reguł
	\item Testowanie algorytmów sztucznej inteligencji (Deep Q-Network) do sterowania
	      agentami
\end{itemize}

\section{Opis Programu}

\subsection{Opis dla Klienta}
System symuluje wybuch epidemii zombie w zamkniętym środowisku. Użytkownik
obserwuje w czasie rzeczywistym:
\begin{itemize}
	\item 5 typów postaci (zombie, ludzie, zarażeni, medycy, żołnierze) reprezentowanych różnymi kolorami
	\item Dynamiczne zmiany populacji w wyniku interakcji między postaciami
	\item Rozprzestrzenianie się infekcji i jej hamowanie poprzez działania medyków i
	      żołnierzy
	\item Licznik czasu rzeczywistego każdego typu postaci
\end{itemize}

\subsection{Opis Szczegółowy}
Symulacja zostaje zainicjalizowana z domyślnymi parametrami:
\begin{itemize}
	\item Rozmiar siatki: 256x256 komórek
	\item Rozdzielczość wyświetlania: 4 piksele na komórkę
	\item Częstotliwość aktualizacji: 100ms między kolejnymi ruchami agentów
	\item Liczba postaci początkowych: 40 zombie, 60 ludzi, 0 zarażonych, 20 medyków, 20
	      żołnierzy
\end{itemize}

Każda postać ma przypisane:
\begin{itemize}
	\item Pozycję (x, y) na siatce
	\item Kolor reprezentujący typ
	\item Logikę comportamentu specyficzną dla typu
	\item Bufor stanu 4 ramek o wymiarach 11x11 dla współpracy z DQN
\end{itemize}

\subsection{Instrukcja Obsługi}

\subsubsection{Uruchomienie Programu}
\begin{enumerate}
	\item Zainstalować zależności: \texttt{pip install -r requirements.txt}
	\item Uruchomić program: \texttt{python main.py}
	\item Program automatycznie inicjalizuje agentów DQN i tworzy siatkę z 40 zombie, 60
	      ludźmi, 20 medykami i 20 żołnierzami
	\item Zamknąć okno aplikacji, aby zakończyć symulację
\end{enumerate}

\subsubsection{Legenda Kolorów}
\begin{itemize}
	\item \textbf{Zielony} - Zombie (agressywne, zarażają ludzi)
	\item \textbf{Niebieski} - Ludzie (defensywni, rozglaszają ostrzeżenia)
	\item \textbf{Pomarańczowy} - Zarażeni (przejściowy stan przed zmianą w zombie)
	\item \textbf{Biały} - Medycy (leczą zarażonych, przywracając ich do poprzedniego stanu)
	\item \textbf{Żółty} - Żołnierze (eliminują zombie w bliskiej okolicy)
\end{itemize}

\subsubsection{Wizualizacja Wyników}
Na górze okna wyświetlany jest licznik aktualnych populacji każdego typu
postaci w formacie: \texttt{Zombies: X | Humans: Y | Infected: Z | Medics: A |
	Soldiers: B}

\subsection{Instrukcja Wdrożenia}

\subsubsection{Wymagania Systemowe}
\begin{itemize}
	\item Python 3.8+
	\item Pygame 2.1.0+
	\item PyTorch 1.10.0+ (dla modułu DQN)
	\item NumPy 1.21.0+
\end{itemize}

\subsubsection{Kroki Wdrażania}
\begin{enumerate}
	\item Klonować/pobrać repozytorium projektu
	\item Utworzyć wirtualne środowisko: \texttt{python -m venv .venv}
	\item Aktywować środowisko: \texttt{source .venv/bin/activate}
	\item Zainstalować pakiety: \texttt{pip install -r requirements.txt}
	\item Uruchomić program: \texttt{python main.py}
\end{enumerate}

\subsection{Dodatkowe Informacje}
\begin{itemize}
	\item Projekt wykorzystuje bibliotekę Pygame do renderowania 2D
	\item DQN (Deep Q-Network) jest zaimplementowany dla potencjalnego uczenia agentów
	\item Testy jednostkowe dostępne w katalogu \texttt{tests/}
	\item Mutacyjne testy dostępne w katalogu \texttt{mutants/}
\end{itemize}

\newpage
% ========== CZĘŚĆ II: MODEL MATEMATYCZNY I ALGORYTMY ==========
\section*{Część II: Model Matematyczny i Algorytmy}

\subsection{Model Matematyczny Systemu}

System symulacyjny oparty jest na modelu wieloagentowym (MAS), w którym agenci
poruszają się po dyskretnej siatce dwuwymiarowej. Decyzje agentów sterowane są
przez algorytm \textit{Deep Q-Network} (DQN), który aproksymuje funkcję
wartości $Q(s, a)$.

\subsubsection{Przestrzeń Obserwacji (Wejście Sieci)}

Dla każdego agenta uczącego się, stan $s_t$ składa się z dwóch heterogenicznych
komponentów, które są przetwarzane równolegle:

\begin{enumerate}
	\item \textbf{Komponent Wizualny (Local Grid Map):}
	      Agent obserwuje lokalne otoczenie w promieniu 5 kratek, co daje macierz o wymiarach $11 \times 11$. Aby zapewnić agentowi rozumienie dynamiki ruchu (kierunek i prędkość obiektów), wejście sieci składa się ze \textbf{stosu 4 kolejnych klatek} (Frame Stacking).
	      \begin{equation}
		      I_t \in \mathbb{R}^{4 \times 11 \times 11}
	      \end{equation}
	      Gdzie wymiar 4 odpowiada kanałom: $[t, t-1, t-2, t-3]$.

	\item \textbf{Komponent Wektorowy (Sensor):}
	      Dodatkowy wektor cech $V_t \in \mathbb{R}^N$, zawierający znormalizowane informacje kierunkowe (np. wektor do najbliższego wroga) oraz stan wewnętrzny (np. przeładowanie broni). Rozmiar $N$ zależy od typu agenta:
	      \begin{itemize}
		      \item \textbf{Zombie ($N=2$):} Wektor do najbliższej ofiary (wykrytej wizualnie lub poprzez sygnał roju \texttt{prey\_spotted}).
		      \item \textbf{Soldier ($N=3$):} Wektor do najbliższego celu (wroga) oraz skalarna wartość gotowości broni ($\delta_{weapon} \in [0, 1]$).
		      \item \textbf{Human ($N=4$):} Wektor złożony z dwóch składowych kierunkowych:
		            \begin{equation}
			            V_{human} = [\vec{v}_{threat}, \vec{v}_{rescue}]
		            \end{equation}
		            Gdzie:
		            \begin{itemize}
			            \item $\vec{v}_{threat} = [dx_t, dy_t]$: Wektor do najbliższego zagrożenia (zombie/infekcja).
			            \item $\vec{v}_{rescue} = [dx_r, dy_r]$: \textbf{Wektor ratunku}, wskazujący na pozycję żołnierza, który nadał sygnał \texttt{zombies\_killed}. Cywil nie kieruje się do dowolnego żołnierza, lecz do tego, który aktywnie eliminuje zagrożenie, traktując jego pozycję jako bezpieczny azyl.
		            \end{itemize}
	      \end{itemize}
\end{enumerate}

\subsubsection{Przestrzeń Akcji ($\mathcal{A}$)}

Przestrzeń akcji jest dyskretna i wynosi $|\mathcal{A}| = 5$:
\begin{equation}
	\mathcal{A} = \{0: \text{Góra}, 1: \text{Dół}, 2: \text{Lewo}, 3: \text{Prawo}, 4: \text{Czekaj}\}
\end{equation}

\subsection{Architektura Sieci Neuronowej (ZombieQNetwork)}

Zastosowano hybrydową architekturę typu \textit{Dual-Stream}, która łączy
przetwarzanie obrazu (CNN) z przetwarzaniem danych numerycznych (MLP). Sieć
składa się z następujących bloków:

\subsubsection{Gałąź Wizualna (CNN)}
Przetwarza stos 4 klatek $11 \times 11$. Składa się z dwóch warstw
konwolucyjnych z funkcją aktywacji ReLU:
\begin{enumerate}
	\item \textbf{Conv2D-1:} Wejście: 4 kanały, Wyjście: 32 filtry, Kernel: $3 \times 3$, Stride: 1, Padding: 1.
	\item \textbf{Conv2D-2:} Wejście: 32 kanały, Wyjście: 64 filtry, Kernel: $3 \times 3$, Stride: 1, Padding: 1.
	\item \textbf{Flatten:} Spłaszczenie mapy cech do wektora o rozmiarze $64 \times 11 \times 11 = 7744$.
\end{enumerate}

\subsubsection{Gałąź Wektorowa (MLP)}
Przetwarza wektor sensorów $V_t$.
\begin{enumerate}
	\item \textbf{Linear:} Transformacja wektora $N$-elementowego do przestrzeni 32 cech.
\end{enumerate}

\subsubsection{Fuzja i Decyzja (Fully Connected)}
Cechy wizualne i wektorowe są łączone (konkatenacja), tworząc wektor o wymiarze
$7744 + 32 = 7776$. Następnie sygnał przechodzi przez warstwy gęste:
\begin{enumerate}
	\item \textbf{FC-1:} Warstwa ukryta, 512 neuronów, aktywacja ReLU.
	\item \textbf{FC-2 (Output):} Warstwa wyjściowa, 5 neuronów (wartości Q dla każdej akcji).
\end{enumerate}

\subsection{Hiperparametry i Proces Uczenia}

Agent uczy się metodą \textit{Experience Replay} z wykorzystaniem dwóch sieci
(Policy Net i Target Net) w celu stabilizacji treningu.

\begin{table}[H]
	\centering
	\caption{Hiperparametry algorytmu DQN}
	\label{tab:dqn_params}
	\begin{tabular}{|l|c|l|}
		\hline
		\textbf{Parametr}        & \textbf{Wartość} & \textbf{Opis}                              \\ \hline
		Input Shape              & $(4, 11, 11)$    & Stos 4 klatek lokalnej mapy                \\ \hline
		Batch Size               & 64               & Liczba próbek w jednym kroku uczenia       \\ \hline
		Learning Rate ($\alpha$) & 0.001            & Współczynnik uczenia optymalizatora Adam   \\ \hline
		Gamma ($\gamma$)         & 0.99             & Współczynnik dyskontowania nagród          \\ \hline
		Replay Buffer            & 10,000           & Pojemność pamięci doświadczeń              \\ \hline
		Target Update            & 1000 kroków      & Częstotliwość aktualizacji sieci docelowej \\ \hline
		Epsilon Start            & 1.0              & Początkowa szansa na losowy ruch           \\ \hline
		Epsilon Decay            & 0.995            & Mnożnik zmniejszający eksplorację          \\ \hline
		Epsilon Min              & 0.05             & Minimalny poziom eksploracji               \\ \hline
	\end{tabular}
\end{table}

\subsection{Kształtowanie Nagrody (Reward Shaping)}

Funkcja nagrody $R(s, a, s')$ została zaprojektowana tak, aby promować
specyficzne zachowania dla każdej klasy.

\subsubsection{Zombie: Agresywna Inteligencja Roju (Swarm Intelligence)}

Zombie działają w oparciu o logikę roju, wykorzystując zarówno wzrok, jak i
komunikację radiową ("krzyk") do lokalizacji ofiar. Kluczowym elementem jest
mechanizm \textit{Oblężenia} (Siege Logic), który eliminuje kary za bezczynność
w tłumie, promując osaczanie ofiary.

\paragraph{Funkcja Nagrody}
Nagroda całkowita $R_{zombie}$ w kroku $t$ jest sumą składowych:
\begin{equation}
	R_{zombie}^{(t)} = R_{infect} + R_{siege} + R_{pursuit} + R_{env}
\end{equation}

Gdzie poszczególne składowe zdefiniowane są następująco:
\begin{itemize}
	\item \textbf{Infekcja ($R_{infect}$):} Dyskretna nagroda za sukces ataku.
	      $$ R_{infect} = \begin{cases} +10 & \text{jeśli ofiara została zarażona} \\ 0 & \text{w p.p.} \end{cases} $$

	\item \textbf{Strefa Oblężenia ($R_{siege}$):} Promuje utrzymywanie pozycji w bezpośrednim sąsiedztwie ofiary ($d_t \le 2.5$), niwelując kary za brak ruchu.
	      $$ R_{siege} = \begin{cases} +0.5 & \text{jeśli } d_t \le 2.5 \\ 0 & \text{w p.p.} \end{cases} $$

	\item \textbf{Pościg ($R_{pursuit}$):} Nagroda za zbliżanie się do celu, aktywna tylko poza strefą oblężenia.
	      $$ R_{pursuit} = \begin{cases}
			      (8.0 - d_t) \cdot 0.2 & \text{jeśli } d_t < 8.0                     \\
			      +0.5                  & \text{jeśli } d_t < d_{t-1} \land d_t > 2.5 \\
			      -0.5                  & \text{jeśli } d_t > d_{t-1} \land d_t > 2.5 \\
			      0                     & \text{w p.p.}
		      \end{cases} $$

	\item \textbf{Kary Środowiskowe ($R_{env}$):}
	      $$ R_{env} = -0.1 (\text{kolizja}) - 0.05 (\text{czas}) + \begin{cases} 0 & \text{jeśli akcja=Wait} \land d_t \le 2.5 \\ -0.5 & \text{jeśli akcja=Wait} \land d_t > 2.5 \end{cases} $$
\end{itemize}

\begin{algorithm}[H]
	\caption{Algorytm Zachowania Zombie (Act)}
	\KwIn{Stan otoczenia, sygnały}

	$target \gets \text{None}, min\_dist \gets \infty$\;

	\tcp{1. Percepcja (Wzrok i Słuch)}
	\For{$char \in \text{VisibleCharacters}$}{
		\If{$dist(me, char) < min\_dist$}{
			$target \gets char$; $min\_dist \gets dist$\;
			\textbf{Broadcast}("prey\_spotted", $char.pos$)\;
		}
	}

	\If{$target$ is None}{
		$target \gets \text{SignalSource("prey\_spotted")}$\;
	}

	\tcp{2. Próba Infekcji}
	\If{$min\_dist \le 1.99$ \textbf{and} $Cooldown == 0$}{
		\textbf{Infect}($target$)\;
		\Return $+10$\;
	}

	\tcp{3. Obliczenie Nagrody za Ruch}
	\If{$min\_dist \le 2.5$}{
		\tcp*[h]{Strefa Oblężenia}
		\Return $+0.5 + R_{base\_prox}$\;
	}
	\Else{
		\tcp*[h]{Strefa Pościgu}
		$R_{dir} \gets (min\_dist < prev\_dist) ? +0.5 : -0.5$\;
		\Return $R_{dir} + R_{base\_prox}$\;
	}
\end{algorithm}

\subsubsection{Soldier: Reaktywna Obrona}

Żołnierz jest jednostką bojową. Jego funkcja celu promuje eliminację zagrożeń oraz taktyczne poruszanie się.

\paragraph{Funkcja Nagrody}
\begin{equation}
	R_{soldier}^{(t)} = R_{kill} + R_{survival} + R_{tactical} + R_{death}
\end{equation}

\begin{itemize}
	\item $R_{kill} = +10.0$: Za skuteczną eliminację zombie.
	\item $R_{survival} = +0.5$: Stała nagroda za każdą klatkę przeżycia.
	\item $R_{death} = -10.0$: Kara terminalna za śmierć.
	\item $R_{tactical}$: Kształtowanie zachowania zależne od stanu broni.
\end{itemize}

\begin{algorithm}[H]
	\caption{Algorytm Zachowania Żołnierza}
	\KwIn{Wzrok, Sygnały radiowe}

	\tcp{1. Wybór Celu (Priorytety)}
	\If{$\exists zombie \in SIGHT\_RANGE$}{
		$target \gets \text{NearestVisibleZombie}$\;
	}
	\ElseIf{$\exists signal \in \text{threat\_alert}$}{
		$target \gets \text{SignalSource}$\;
	}
	\tcp{2. Walka}
	\If{$WeaponCooldown == 0$ \textbf{and} $dist(target) \le 3.0$}{
		\textbf{Shoot}($target$)\;
		\textbf{Broadcast}("zombies\_killed")\;
		\Return $+10.0$\;
	}

	\tcp{3. Ocena Ruchu}
	\Return $R_{survival} + R_{tactical}(dist\_delta, weapon\_status)$\;
\end{algorithm}

\subsubsection{Human: Instynkt Przetrwania i Ewhakuacja}

Model cywila jest najbardziej złożony pod względem motywacji, balansując między
panicznym strachem (ucieczka od zagrożenia) a nadzieją (dążenie do strefy, w
której słychać strzały żołnierzy). Sieć neuronowa otrzymuje na wejściu
hybrydowy wektor sterujący, łączący te dwie sprzeczne siły.

\paragraph{Funkcja Nagrody}
Całkowita nagroda $R_{human}^{(t)}$ jest sumą składowych promujących
przetrwanie i karzących za ryzykowne zachowania:

\begin{equation}
	R_{human}^{(t)} = R_{survival} + R_{escape} + R_{panic} + R_{collision} + R_{death}
\end{equation}

Szczegółowa definicja składowych:

\begin{itemize}
	\item \textbf{Przetrwanie ($R_{survival}$):} Stała nagroda za każdą klatkę pozostania przy życiu.
	      $$ R_{survival} = +0.1 $$

	\item \textbf{Gradient Ucieczki ($R_{escape}$):} Nagroda dynamiczna, zależna od zmiany dystansu do najbliższego zagrożenia ($d_t$).
	      $$ R_{escape} = \begin{cases}
			      +0.5 & \text{jeśli } d_t > d_{t-1} \text{ (Skuteczna ucieczka)}       \\
			      -0.5 & \text{jeśli } d_t < d_{t-1} \text{ (Bieg w stronę zagrożenia)} \\
			      0    & \text{w p.p.}
		      \end{cases} $$

	\item \textbf{Kara Paniczna ($R_{panic}$):} Funkcja aktywująca się po przekroczeniu progu bezpieczeństwa $D_{safe} = 5.0$. Kara rośnie liniowo wraz ze zbliżaniem się do zombie.
	      $$ R_{panic} = -\max(0, (5.0 - d_t) \cdot 0.2) $$

	\item \textbf{Kary Środowiskowe ($R_{collision}$):} Penalizacja za uderzenie w przeszkodę (ścianę), co zapobiega blokowaniu się agenta w rogach mapy.
	      $$ R_{collision} = \begin{cases} -0.5 & \text{w przypadku kolizji} \\ 0 & \text{w p.p.} \end{cases} $$

	\item \textbf{Śmierć ($R_{death}$):} Terminalna kara za zostanie zarażonym.
	      $$ R_{death} = -10.0 $$
\end{itemize}

\begin{algorithm}[H]
	\SetAlgoLined
	\KwData{Pozycja własna, Grid, Sygnały radiowe, $d_{prev}$}
	\KwResult{Wartość nagrody $R_t$}

	$R_t \leftarrow 0.1$ \tcp*{Survival Reward}
	$d_{min} \leftarrow \infty$; $target_{threat} \leftarrow$ None; $target_{rescue} \leftarrow$ None\;

	\tcp{1. Percepcja (Wzrok i Słuch)}
	\For{każda postać $c$ w zasięgu percepcji}{
		\If{$c$ jest Zombie/Infected}{
			$d \leftarrow dist(self, c)$\;
			\If{$d < d_{min}$}{ $d_{min} \leftarrow d$; $target_{threat} \leftarrow c$ }
		}
	}

	\tcp{2. Komunikacja (Broadcast)}
	\If{$d_{min} \le DetectionRange$}{
		\textbf{Broadcast}("threat\_alert", $target_{threat}.pos$)\;
	}

	\tcp{3. Fuzja Sensorów (Wejście Sieci)}
	\If{$target_{threat}$ exists}{
		$\vec{v}_{fear} \leftarrow \text{Normalize}(pos_{self} - pos_{threat})$ \tcp*{Wektor OD zagrożenia}
	}
	\Else{
		$\vec{v}_{fear} \leftarrow [0, 0]$\;
	}

	$S_{rescue} \leftarrow \text{GetSignal("zombies\_killed")}$\;
	\If{$S_{rescue}$ exists}{
		$\vec{v}_{hope} \leftarrow \text{Normalize}(pos_{signal} - pos_{self})$ \tcp*{Wektor DO ratunku}
	}
	\Else{
		$\vec{v}_{hope} \leftarrow [0, 0]$\;
	}

	\textbf{NeuralInput} $\leftarrow [\vec{v}_{fear}, \vec{v}_{hope}]$ \tcp*{Wektor 4-elementowy}

	\tcp{4. Obliczenie Nagrody}
	\If{kolizja ze ścianą}{
		$R_t \leftarrow R_t - 0.5$\;
	}

	\If{$d_{min} < \infty$}{
		\tcp{Strefa Paniki}
		\If{$d_{min} < 5.0$}{
			$R_t \leftarrow R_t - (5.0 - d_{min}) \cdot 0.2$\;
		}

		\tcp{Ocena kierunku ruchu}
		\eIf{$d_{min} > d_{prev}$}{
			$R_t \leftarrow R_t + 0.5$ \tcp*{Ucieczka}
		}{
			$R_t \leftarrow R_t - 0.5$ \tcp*{Zły kierunek}
		}
	}

	\caption{Algorytm Zachowania Człowieka (Human Agent)}
\end{algorithm}

\subsection{Algorytmy Deterministyczne (Medyk)}

Medycy w systemie nie używają sieci neuronowych, lecz hybrydowego podejścia
opartego na algorytmie A* z dynamiczną koordynacją (protokół Claim \& Yield)
oraz ograniczonym polem widzenia ("Fog of War").
\begin{equation}
	\mathcal{O}_{medic} = \{ p_j \mid p_j \in \text{Agents} \land p_j \ne \text{Target} \land \|p_{med} - p_j\| \le R_{vision} \}
\end{equation}
Medyk dodaje agentów do zbioru przeszkód $\mathcal{O}$ tylko wtedy, gdy znajdują się w jego promieniu wzroku ($R_{vision} = 8.0$), co symuluje brak pełnej wiedzy o mapie.

\section{Implementacja Systemu}

\subsection{Architektura Projektu}

Projekt jest podzielony na moduły funkcjonalne:

\begin{itemize}
	\item \textbf{character.py} - Klasa bazowa dla wszystkich agentów
	\item \textbf{zombie.py} - Implementacja zombie (agenci atakujący)
	\item \textbf{human.py} - Implementacja ludzi (agenci defensywni)
	\item \textbf{infected.py} - Implementacja zarażonych (stan przejściowy)
	\item \textbf{medic.py} - Implementacja medyków (agenci leczący)
	\item \textbf{soldier.py} - Implementacja żołnierzy (agenci walczący)
	\item \textbf{grid.py} - Menedżer siatki i wszystkich agentów
	\item \textbf{constants.py} - Stałe konfiguracyjne systemu
	\item \textbf{main.py} - Główna pętla gry z inicjalizacją agentów DQN
	\item \textbf{DQN/learn.py} - Implementacja agentów DQN
	\item \textbf{DQN/q\_network.py} - Definicja sieci neuronowej Q-Network
\end{itemize}

\subsection{Klasa Bazowa Character}

\subsubsection{Opis}
Klasa \texttt{Character} jest abstrakcyjną klasą bazową dla wszystkich typów
agentów. Definiuje wspólne zachowania i atrybuty.

\subsubsection{Atrybuty}
\begin{itemize}
	\item \texttt{x, y} - Aktualna pozycja na siatce
	\item \texttt{grid} - Referencja do obiektu siatki
	\item \texttt{char\_type\_name} - Nazwa typu postaci
	\item \texttt{color} - Kolor reprezentujący typ (RGB)
	\item \texttt{move\_speed} - Czas między ruchami w ms
	\item \texttt{state\_buffer} - Deque przechowujący ostatnie 4 obserwacje
	\item \texttt{signals} - Lista odebranych sygnałów
\end{itemize}

\subsubsection{Metody Kluczowe}

\textbf{update(current\_time)}
Metoda wywoływana każdą klatkę. Sprawdza, czy upłynął czas wymagany do ruchu, i jeśli tak, wywołuje \texttt{move()}, \texttt{act()} i \texttt{process\_signals()}.

\begin{lstlisting}[language=Python]
def update(self, current_time):
    if current_time - self.last_move_time >= self.move_speed:
        self.move()
        self.act()
        self.process_signals()
        self.last_move_time = current_time
	\end{lstlisting}

\textbf{move(action\_code)}
Przemieszcza agenta o jeden krok w kierunku określonym przez \texttt{action\_code}. Sprawdza kolizje i granice siatki.

\begin{lstlisting}[language=Python]
def move(self, action_code=None):
    dx, dy = 0, 0
    if action_code == 0: dy = -1      # Up
    elif action_code == 1: dy = 1     # Down
    elif action_code == 2: dx = -1    # Left
    elif action_code == 3: dx = 1     # Right
    
    new_x = max(0, min(GRID_SIZE - 1, self.x + dx))
    new_y = max(0, min(GRID_SIZE - 1, self.y + dy))
    
    # Check collision with other characters
    occupied = False
    if self.grid:
        for character in self.grid.characters:
            if character is not self and \
               character.x == new_x and \
               character.y == new_y:
                occupied = True
                break
    
    if not occupied:
        self.x = new_x
        self.y = new_y
        return True
    return False
	\end{lstlisting}

\textbf{get\_observation()}
Generuje macierz obserwacji lokalnego otoczenia o wymiarach 11x11.

\begin{lstlisting}[language=Python]
def get_observation(self):
    """Generate 11x11 observation around character"""
    obs = np.zeros((11, 11))
    view_range = 5  # 5 cells in each direction
    
    for dx in range(-view_range, view_range + 1):
        for dy in range(-view_range, view_range + 1):
            world_x = self.x + dx
            world_y = self.y + dy
            
            if 0 <= world_x < GRID_SIZE and \
               0 <= world_y < GRID_SIZE:
                # Check what's at this position
                local_x = dx + view_range
                local_y = dy + view_range
                obs[local_y, local_x] = \
                    self.grid.global_map[world_y, world_x]
    
    return obs
	\end{lstlisting}

\subsection{Klasa Zombie}

\subsubsection{Opis}
Zombie są głównymi antagonistami. Aktywnie szukają i zarażają ludzi, otrzymując
nagrody za zbliżanie się i zarażanie.

\subsubsection{Zachowanie act()}
\begin{lstlisting}[language=Python]
def act(self):
    """Perform aggressive behavior - infect nearby humans"""
    if not self.grid:
        return 0
    
    current_time = pygame.time.get_ticks()
    step_reward = 0
    
    # Find closest human
    min_dist = float('inf')
    for char in self.grid.characters:
        if isinstance(char, Human):
            dist = math.sqrt((self.x - char.x)**2 + 
                           (self.y - char.y)**2)
            if dist < min_dist:
                min_dist = dist
    
    # Reward shaping: approach humans
    if min_dist < float('inf'):
        if min_dist < 5.0:  # VIEW_RANGE
            step_reward += (5.0 - min_dist) * 0.2
            
            if min_dist < self.prev_dist:
                step_reward += 0.5
            elif min_dist > self.prev_dist:
                step_reward -= 0.5
        
        self.prev_dist = min_dist
    
    # Infection logic
    if current_time - self.last_infection_time >= \
       self.INFECTION_COOLDOWN:
        for character in self.grid.characters:
            if isinstance(character, (Human, Medic, Soldier)):
                distance = math.sqrt(
                    (self.x - character.x)**2 + 
                    (self.y - character.y)**2)
                
                if distance <= self.INFECTION_RANGE:
                    self.infect_character(character)
                    step_reward += 10
                    self.last_infection_time = current_time
                    break
    
    return step_reward
	\end{lstlisting}

\subsection{Klasa Human}

\subsubsection{Opis}
Ludzie są defensywni. Unikają zombie, rozglaszają ostrzeżenia i współpracują z
innymi postaciami.

\subsubsection{Zachowanie act()}
Ludzie skanują otoczenie w poszukiwaniu zombie/zarażonych. Jeśli zostaną
wykryte, rozglaszają informację o zagrożeniu w promieniu 7.0 komórek.

\subsection{Klasa Grid}

\subsubsection{Opis}
Klasa \texttt{Grid} zarządza całą symulacją. Przechowuje wszystkich agentów,
obsługuje aktualizacje i renderowanie.

\subsubsection{Metody Kluczowe}

\textbf{get\_global\_map\_matrix()}
Tworzy macierz 256x256 reprezentującą stan całej siatki.

\begin{lstlisting}[language=Python]
def get_global_map_matrix(self):
    """Create global map matrix"""
    matrix = np.zeros((GRID_SIZE, GRID_SIZE))
    
    for char in self.characters:
        if isinstance(char, (Zombie, Infected)):
            matrix[char.y, char.x] = -1.0  # Enemies
        elif isinstance(char, (Human, Medic, Soldier)):
            matrix[char.y, char.x] = 0.5   # Allies
    
    return matrix
	\end{lstlisting}

\textbf{get\_stats()}
Zwraca liczbę każdego typu postaci.

\subsection{Główna Pętla Gry}

\begin{lstlisting}[language=Python]
def main():
    """Main game loop"""
    print("Initializing simulation...")
    screen = pygame.display.set_mode((WINDOW_SIZE, WINDOW_SIZE+ 40))
    pygame.display.set_caption("Zombie Outbreak Simulation")
    clock = pygame.time.Clock()
    
    # Initialize DQN agents
    zombie_agent = DQNAgent(input_shape=(4, 11, 11))
    human_agent = DQNAgent(input_shape=(4, 11, 11), vector_size=4)
    soldier_agent = DQNAgent(input_shape=(4, 11, 11), vector_size=3)
    
    grid = Grid(num_zombies=40, num_humans=60, num_infected=0, num_medics=20, num_soldiers=20)
    font = pygame.font.Font(None, 24)
    
    running = True
    while running:
        for event in pygame.event.get():
            if event.type == pygame.QUIT:
                running = False
        
        # Update game logic
        global_map = grid.get_global_map_matrix()
        update_game_logic(grid, global_map)
        
        # Render
        screen.fill(COLOR_BACKGROUND)
        grid.draw(screen)
        
        # Display stats
        stats = grid.get_stats()
        stats_text = f"Zombies: {stats[0]} | Humans: {stats[1]}"
        stats_surface = font.render(stats_text, True, 
                                   (255, 255, 255))
        screen.blit(stats_surface, (10, 10))
        
        pygame.display.flip()
        clock.tick(FPS)
    
    pygame.quit()
	\end{lstlisting}

\newpage

\section{Testy}

\subsection{Testy Jednostkowe}

\subsubsection{Zakres Testów}

Projekt zawiera kompleksowe testy jednostkowe w pliku
\texttt{tests/character\_test.py} obejmujące:
\begin{enumerate}
	\item \textbf{Inicjalizacja} - Sprawdzenie poprawnego utworzenia obiektów
	\item \textbf{Ruch} - Weryfikacja logiki poruszania się
	\item \textbf{Obserwacja} - Testowanie generowania macierzy obserwacji
	\item \textbf{Sygnały} - Testowanie systemu komunikacji
	\item \textbf{Kolizje} - Sprawdzenie unikania kolizji
\end{enumerate}

\subsubsection{Przykładowe Testy}

\textbf{Test inicjalizacji stanu bufora}
\begin{lstlisting}[language=Python]
def test_character_state_buffer_initialization(character):
    """Test state buffer is initialized with 4 frames"""
    assert len(character.state_buffer) == 4
    for frame in character.state_buffer:
        assert frame.shape == (11, 11)
        assert np.all(frame == 0)
	\end{lstlisting}

\textbf{Test ruchu w górę}
\begin{lstlisting}[language=Python]
def test_move_up(character):
    """Test move action_code 0 moves up (dy = -1)"""
    initial_x, initial_y = character.x, character.y
    character.move(action_code=0)
    assert character.x == initial_x
    assert character.y == initial_y - 1
	\end{lstlisting}

\subsubsection{Wyniki Testów}

Polecenie: \texttt{pytest tests/character\_test.py -v}

\textbf{Statystyka testów:}
\begin{itemize}
	\item Liczba testów: 63
	\item Testy przechodzące: 63
	\item Pokrycie kodu: ~100\% klasy Character
\end{itemize}

\subsection{Testy Mutacyjne}

\subsubsection{Cel Testów Mutacyjnych}

Testy mutacyjne sprawdzają, czy zmiany w kodzie są wykrywane przez testy
jednostkowe. System MutMut automatycznie wprowadza małe zmiany (mutacje) w
kodzie i sprawdza, czy testy je wykrywają (killed) czy nie (survived).

\subsubsection{Strategia Testowania}

\begin{enumerate}
	\item Uruchomienie: \texttt{mutmut run}
	\item Analiza wyników: \texttt{mutmut results}
	\item Wyświetlenie wszystkich mutacji, które przeżyły: \texttt{mutmut results}
	\item Wyświetlenie wszystkich mutacji (pełna lista): \texttt{mutmut results
		      --all=true}
	\item Wyświetlenie szczegółów: \texttt{\detokenize{mutmut show
			      character.x|Character|get_target_vector__mutmut_6}}
\end{enumerate}

\subsubsection{Przykładowe Mutacje}

\begin{itemize}
	\item \textbf{Mutacja 22} - Zmiana operatora \texttt{>=} na \texttt{>} w warunku ruchu
	      \begin{lstlisting}[language=Python]
# Original
if current_time - self.last_move_time >= self.move_speed:
# Mutated
if current_time - self.last_move_time > self.move_speed:
		\end{lstlisting}
	      Status: \textbf{KILLED} (testy wykrywają tę zmianę)
	\item \textbf{Mutacja 23} - Zmiana wartości zwracanej przez funkcję \texttt{move()} z \texttt{True} na \texttt{False}
	      \begin{lstlisting}[language=Python]
# Original
if not occupied:
	self.x = new_x
	self.y = new_y
	return True
# Mutated
if not occupied:
	self.x = new_x
	self.y = new_y
	return False
		  \end{lstlisting}
	      Status: \textbf{KILLED} (testy wykrywają tę zmianę)
\end{itemize}
\subsubsection{Wyniki Testów Mutacyjnych}
\begin{itemize}
	\item Liczba mutacji: 133
	\item Mutacje zabite (killed): 133
	\item Mutacje przeżyłe (survived): 0
	\item Wskaźnik zabicia mutacji: $\frac{133}{133} \times 100\% = 100\%$
\end{itemize}

\newpage

\section{Diagram UML}

\subsection{Hierarchia Klas}

\textbf{Struktura dziedziczenia:}

\begin{verbatim}
	Character (abstract)
	  - Zombie
	  - Human
	    - Medic
	    - Soldier
	  - Infected
	
	Grid
	  - characters: List[Character]
	\end{verbatim}

\subsection{Diagram Klas Głównych}

\begin{verbatim}
	+-Character---------+
	| - x: int          |
	| - y: int          |
	| - grid: Grid      |
	| - state_buffer    |
	| - color: RGB      |
	|+move()            |
	|+act()             |
	|+get_observation() |
	|+draw(screen)      |
	+-------------------+
	         ^
	         |
	    +----+-----+-------+
	    |          |       |
	  Zombie    Human   Medic  Soldier
	           |
	        Infected
	\end{verbatim}

\newpage

\section{Analiza Złożoności Obliczeniowej}

\subsection{Złożoność Czasowa}

\subsubsection{Operacje na Agentach}

Główne operacje wykonywane na każdego agenta (dla każdej klatki):

\begin{itemize}
	\item \textbf{Ruch:} $O(n)$ - sprawdzenie kolizji z wszystkimi agentami
	\item \textbf{Obserwacja:} $O(1)$ - generowanie matrycy 11x11 (stały rozmiar)
	\item \textbf{Działanie (act):} $O(n)$ - skanowanie w poszukiwaniu celów/zagrożeń
	\item \textbf{Komunikacja:} $O(n)$ - rozsyłanie sygnałów do agentów w zasięgu
\end{itemize}

gdzie $n$ = liczba agentów w systemie

\subsubsection{Pętla Główna Gry}

Każda klatka wykonuje:
\begin{align}
	O_{frame} & = O(n \cdot O_{agent}) \\
	          & = O(n) \cdot O(n)      \\
	          & = O(n^2)
\end{align}

Dla domyślnej konfiguracji: - $n = 100$ agentów (40 zombie + 60 ludzi) -
$O_{frame} = O(100^2) = 10,000$ operacji podstawowych na klatkę - FPS = 60 →
$60 \times 10,000 = 600,000$ operacji/sekundę

\textbf{Wniosek:} Złożoność $O(n^2)$ jest akceptowalna dla zakresu liczby agentów (do ~200).

\subsection{Złożoność Pamięciowa}

\subsubsection{Pamięć na Agenta}

\begin{itemize}
	\item \textbf{Atrybuty proste:} $x, y, color, move\_speed$ = 4 × 32-bit = 16 bajtów
	\item \textbf{Referencja do gridu:} 1 wskaźnik = 8 bajtów
	\item \textbf{State buffer:} 4 × (11 × 11) floats = 4 × 121 × 8 = 3,872 bajtów
	\item \textbf{Lista sygnałów:} średnio 5 sygnałów × 32 bajty = 160 bajtów
	\item \textbf{Inne atrybuty:} ~100 bajtów
\end{itemize}

\textbf{Razem na agenta:} $\approx 4,200$ bajtów $\approx 4$ KB

\subsubsection{Całkowita Pamięć Systemu}

$$M_{total} = n \cdot 4\text{ KB} + M_{grid} + M_{pygame}$$

Dla $n=100$: $M_{total} \approx 100 \times 4\text{ KB} + 500\text{ KB} +
	2\text{ MB} \approx 2.9\text{ MB}$

\textbf{Wniosek:} Footprint pamięciowy jest znikomy dla współczesnych komputerów.

\subsection{Optymalizacje}

\subsubsection{Spatial Hashing}

Obecnie $O(n^2)$ jest nieuniknione ze względu na potrzebę sprawdzania
wszystkich agentów. Możliwa optymalizacja:

\textbf{Spatial Grid Hash Map:}
\begin{itemize}
	\item Podzielić siatkę na komórki (np. 16×16)
	\item Każdy agent przechowywać w komórce, w której się znajduje
	\item Podczas sprawdzania kolizji, sprawdzać tylko agentów w sąsiadujących komórkach
	\item Złożoność ulepszona do: $O(n \cdot k)$ gdzie $k$ = średnia liczba agentów w
	      sąsiadujących komórkach ($k \ll n$)
\end{itemize}

\subsubsection{Lazy Evaluation}

\begin{itemize}
	\item Macierz obserwacji generować tylko gdy będzie potrzebna (dla agentów z DQN)
	\item Sygnały buforować i przetwarzać grupowo
\end{itemize}

\subsection{Asymptotyczna Analiza}

\begin{table}[H]
	\centering
	\begin{tabular}{|c|c|c|c|}
		\hline
		\textbf{Operacja} & \textbf{Złożoność} & \textbf{Dla n=100} & \textbf{Dla n=1000} \\
		\hline
		Ruch agenta       & $O(n)$             & 100 ops            & 1,000 ops           \\
		Act (szukanie)    & $O(n)$             & 100 ops            & 1,000 ops           \\
		Pełna klatka      & $O(n^2)$           & 10,000 ops         & 1,000,000 ops       \\
		\hline
	\end{tabular}
	\caption{Analiza złożoności dla różnych rozmiarów systemu}
\end{table}

\newpage

\section{Wyniki Eksperymentów i Testy Porównawcze}

\subsection{Testy Wydajnościowe}

\subsubsection{Test 1: Stabilność FPS}

\textbf{Konfiguracja:} 100 agentów (40Z, 60H), 60 FPS, czas trwania 1 minuta

\textbf{Rezultaty:}
\begin{itemize}
	\item Średni FPS: 59.8
	\item Min FPS: 58
	\item Max FPS: 61
	\item Stabilność: Bardzo wysoka
\end{itemize}

\textbf{Wniosek:} System utrzymuje stałą prędkość dla domyślnej konfiguracji.

\subsubsection{Test 2: Skalowanie z Liczbą Agentów}

\begin{table}[H]
	\centering
	\begin{tabular}{|c|c|c|c|}
		\hline
		\textbf{Liczba Agentów} & \textbf{Średni FPS} & \textbf{Czas/Frame (ms)} & \textbf{CPU \%} \\
		\hline
		50                      & 60                  & 16.7                     & 15\%            \\
		100                     & 59.8                & 16.7                     & 28\%            \\
		200                     & 58                  & 17.2                     & 52\%            \\
		500                     & 35                  & 28.6                     & 95\%            \\
		\hline
	\end{tabular}
	\caption{Wpływ liczby agentów na wydajność}
\end{table}

\textbf{Obserwacja:} Wydajność degrada się kwadratowo, potwierdzając analizę złożoności $O(n^2)$.

\subsection{Test Dynamiki Populacji}

\subsubsection{Scenariusz: 40 Zombie vs 60 Ludzi}

\textbf{Hipoteza:} Zombie powoli zarazą większość ludzi, tworząc wachlarzy epidemii.

\textbf{Obserwacja:}
\begin{itemize}
	\item Pierwsze zarażenie: ~10 sekund
	\item Pik zarażonych: ~60 sekund (30\% populacji)
	\item Równowaga: ~180 sekund (50Z, 30H, 20Inf)
	\item Trend: Exponential growth zarażonych, linear decay ludzi
\end{itemize}

\subsubsection{Test z Medykami}

\textbf{Konfiguracja:} 40Z, 50H, 5 Medyków

\textbf{Wynik:} Medycy znacznie opóźniają rozprzestrzenianie się epidemii, choć asymptotycznie nie są w stanie jej powstrzymać.

\subsection{Analiza Statystyczna}

\subsubsection{Współczynnik Infekacji}

Współczynnik $R_0$ (reproduction number) - ile osób zaraża jeden zombie:

$$R_0 = \frac{\text{liczba zarażonych w równowadze}}{\text{liczba zombie} \cdot \text{czas}}$$

Empirycznie: $R_0 \approx 1.5$ dla domyślnych parametrów.

\subsubsection{Czas Połowicznego Rozpadu Ludzi}

$$t_{1/2} = \frac{\ln(2) \cdot N_0}{k}$$

gdzie $k$ = stała rozpadu (zarażania)

Empirycznie: $t_{1/2} \approx 120$ sekund dla 100 agentów.

\newpage

\section{Wybrane Fragmenty Kodu}

\subsection{Fragment 1: Logika Zarażenia Zombie}

\begin{lstlisting}[language=Python]
# zombie.py
def infect_character(self, character):
    """Convert a character to infected"""
    if self.grid:
        from infected import Infected
        
        idx = self.grid.characters.index(character)
        infected = Infected(
            character.x, character.y, self.grid, 
            previous_type=character.get_type_name()
        )
        self.grid.characters[idx] = infected
    
    character.get_infected()
    return True
	\end{lstlisting}

\textbf{Opis:} Metoda zastępuje postawę ludzie zainfekowanym obiektem. Zapisuje poprzedni typ dla możliwości leczenia.

\subsection{Fragment 2: Generowanie Macierzy Obserwacji}

\begin{lstlisting}[language=Python]
# character.py
def get_observation(self):
    """Generate 11x11 observation matrix"""
    obs = np.zeros((11, 11))
    view_range = 5
    
    for dx in range(-view_range, view_range + 1):
        for dy in range(-view_range, view_range + 1):
            world_x = self.x + dx
            world_y = self.y + dy
            
            if 0 <= world_x < GRID_SIZE and \
               0 <= world_y < GRID_SIZE:
                local_x = dx + view_range
                local_y = dy + view_range
                obs[local_y, local_x] = \
                    self.grid.global_map[world_y, world_x]
    
    self.state_buffer.append(obs)
    return np.array(list(self.state_buffer))
	\end{lstlisting}

\textbf{Opis:} Tworzy macierz 11x11 reprezentującą otoczenie agenta. Skanuje grid wokół pozycji i mapuje wartości z globalnej mapy.

\subsection{Fragment 3: System Nagród DQN}

\begin{lstlisting}[language=Python]
# zombie.py - act() method reward shaping
step_reward = 0
if min_dist < float('inf'):
    if min_dist < 5.0:  # VIEW_RANGE
        # Reward for approaching humans
        step_reward += (5.0 - min_dist) * 0.2
        
        # Reward shaping for correct direction
        if min_dist < self.prev_dist:
            step_reward += 0.5  # Getting closer
        elif min_dist > self.prev_dist:
            step_reward -= 0.5  # Getting farther
    
    self.prev_dist = min_dist

# Infection reward
if distance <= self.INFECTION_RANGE:
    self.infect_character(character)
    step_reward += 10  # Major reward for infection
	\end{lstlisting}

\textbf{Opis:} System kształtowania nagród (reward shaping) dla zombie. Rozdaje nagrody za zbliżanie się do ludzi i karę za oddalanie się. Największa nagroda za rzeczywiste zarażenie.

\subsection{Fragment 4: Obsługa Sygnałów}

\begin{lstlisting}[language=Python]
# human.py
def broadcast_threat_info(self, threats):
    """Broadcast threat information to nearby characters"""
    if not self.grid:
        return
    
    for character in self.grid.characters:
        if character is self:
            continue
        
        dx = self.x - character.x
        dy = self.y - character.y
        distance = math.sqrt(dx*dx + dy*dy)
        
        if distance <= self.THREAT_BROADCAST_RANGE:
            signal = {
                'from': self,
                'threats': threats,
                'timestamp': pygame.time.get_ticks()
            }
            character.signals.append(signal)
	\end{lstlisting}

\textbf{Opis:} Człowiek rozgłasza informacje o zagrożeniu wszystkim w zasięgu. Każdy agent otrzymuje sygnał zawierający rodzaj zagrożenia i lokalizację.

\newpage

\section{Podsumowanie i Wnioski}

\subsection{Osiągnięte Cele}

\begin{enumerate}
	\item [\checkmark] Zaimplementowano wieloagentowy system symulacji epidemiologicznej
	\item [\checkmark] Pięć typów agentów z różnymi zachowaniami i celami
	\item [\checkmark] System komunikacji i rozprzestrzeniania informacji
	\item [\checkmark] Mapa globalnych obserwacji dla potencjalnego uczenia DQN
	\item [\checkmark] Kompleksowe testy jednostkowe i mutacyjne
	\item [\checkmark] Interfejs graficzny w Pygame
\end{enumerate}

\subsection{Zaobserwowane Zachowania Emergentne}

\begin{itemize}
	\item Zombie naturalnie formują "swarmy" w kierunku gromad ludzi
	\item Ludzie tworzą grupy ochronne wokół medyków
	\item Oszacowana dynamika populacji odpowiada modelom epidemiologicznym (SIR)
	\item Żołnierze efektywnie redukują populację zombie w lokalnych regionach
\end{itemize}

\subsection{Możliwe Rozszerzenia}

\begin{enumerate}
	\item \textbf{Pełne uczenie DQN} - Wytrenować sieci neuronowe do optymalnego sterowania agentami
	\item \textbf{Spatial Hashing} - Optymalizacja do $O(n \log n)$ dla większych systemów
	\item \textbf{Różne środowiska} - Przeszkody, budynki, strefy bezpieczne
	\item \textbf{Genetyczne Algorytmy} - Ewolucja parametrów behawioralnych agentów
	\item \textbf{Analiza Sieci Społecznej} - Modelowanie połączeń między agentami
	\item \textbf{Różne patogeny} - Zmienność infekcji, mutacje zombie
\end{enumerate}

\subsection{Wnioski Techniczne}

\begin{itemize}
	\item Architektura oparta na dziedziczeniu (Character) pozwala na łatwe dodawanie
	      nowych typów agentów
	\item System buforowania obserwacji (state\_buffer) jest efektywny dla DQN
	\item Reward shaping jest kluczowy dla osłabiania podejmowania decyzji przez agentów
	\item Kompleksowe testy (jednostkowe, mutacyjne) zapewniają niezawodność kodu
	\item Złożoność $O(n^2)$ jest akceptowalna do ~200 agentów
\end{itemize}

\end{document}
